% Options for packages loaded elsewhere
\PassOptionsToPackage{unicode}{hyperref}
\PassOptionsToPackage{hyphens}{url}
\documentclass[
  ignorenonframetext,
]{beamer}
\newif\ifbibliography
\usepackage{pgfpages}
\setbeamertemplate{caption}[numbered]
\setbeamertemplate{caption label separator}{: }
\setbeamercolor{caption name}{fg=normal text.fg}
\beamertemplatenavigationsymbolsempty
% remove section numbering
\setbeamertemplate{part page}{
  \centering
  \begin{beamercolorbox}[sep=16pt,center]{part title}
    \usebeamerfont{part title}\insertpart\par
  \end{beamercolorbox}
}
\setbeamertemplate{section page}{
  \centering
  \begin{beamercolorbox}[sep=12pt,center]{section title}
    \usebeamerfont{section title}\insertsection\par
  \end{beamercolorbox}
}
\setbeamertemplate{subsection page}{
  \centering
  \begin{beamercolorbox}[sep=8pt,center]{subsection title}
    \usebeamerfont{subsection title}\insertsubsection\par
  \end{beamercolorbox}
}
% Prevent slide breaks in the middle of a paragraph
\widowpenalties 1 10000
\raggedbottom
\AtBeginPart{
  \frame{\partpage}
}
\AtBeginSection{
  \ifbibliography
  \else
    \frame{\sectionpage}
  \fi
}
\AtBeginSubsection{
  \frame{\subsectionpage}
}
\usepackage{iftex}
\ifPDFTeX
  \usepackage[T1]{fontenc}
  \usepackage[utf8]{inputenc}
  \usepackage{textcomp} % provide euro and other symbols
\else % if luatex or xetex
  \usepackage{unicode-math} % this also loads fontspec
  \defaultfontfeatures{Scale=MatchLowercase}
  \defaultfontfeatures[\rmfamily]{Ligatures=TeX,Scale=1}
\fi
\usepackage{lmodern}
\ifPDFTeX\else
  % xetex/luatex font selection
\fi
% Use upquote if available, for straight quotes in verbatim environments
\IfFileExists{upquote.sty}{\usepackage{upquote}}{}
\IfFileExists{microtype.sty}{% use microtype if available
  \usepackage[]{microtype}
  \UseMicrotypeSet[protrusion]{basicmath} % disable protrusion for tt fonts
}{}
\makeatletter
\@ifundefined{KOMAClassName}{% if non-KOMA class
  \IfFileExists{parskip.sty}{%
    \usepackage{parskip}
  }{% else
    \setlength{\parindent}{0pt}
    \setlength{\parskip}{6pt plus 2pt minus 1pt}}
}{% if KOMA class
  \KOMAoptions{parskip=half}}
\makeatother
\usepackage{color}
\usepackage{fancyvrb}
\newcommand{\VerbBar}{|}
\newcommand{\VERB}{\Verb[commandchars=\\\{\}]}
\DefineVerbatimEnvironment{Highlighting}{Verbatim}{commandchars=\\\{\}}
% Add ',fontsize=\small' for more characters per line
\usepackage{framed}
\definecolor{shadecolor}{RGB}{248,248,248}
\newenvironment{Shaded}{\begin{snugshade}}{\end{snugshade}}
\newcommand{\AlertTok}[1]{\textcolor[rgb]{0.94,0.16,0.16}{#1}}
\newcommand{\AnnotationTok}[1]{\textcolor[rgb]{0.56,0.35,0.01}{\textbf{\textit{#1}}}}
\newcommand{\AttributeTok}[1]{\textcolor[rgb]{0.13,0.29,0.53}{#1}}
\newcommand{\BaseNTok}[1]{\textcolor[rgb]{0.00,0.00,0.81}{#1}}
\newcommand{\BuiltInTok}[1]{#1}
\newcommand{\CharTok}[1]{\textcolor[rgb]{0.31,0.60,0.02}{#1}}
\newcommand{\CommentTok}[1]{\textcolor[rgb]{0.56,0.35,0.01}{\textit{#1}}}
\newcommand{\CommentVarTok}[1]{\textcolor[rgb]{0.56,0.35,0.01}{\textbf{\textit{#1}}}}
\newcommand{\ConstantTok}[1]{\textcolor[rgb]{0.56,0.35,0.01}{#1}}
\newcommand{\ControlFlowTok}[1]{\textcolor[rgb]{0.13,0.29,0.53}{\textbf{#1}}}
\newcommand{\DataTypeTok}[1]{\textcolor[rgb]{0.13,0.29,0.53}{#1}}
\newcommand{\DecValTok}[1]{\textcolor[rgb]{0.00,0.00,0.81}{#1}}
\newcommand{\DocumentationTok}[1]{\textcolor[rgb]{0.56,0.35,0.01}{\textbf{\textit{#1}}}}
\newcommand{\ErrorTok}[1]{\textcolor[rgb]{0.64,0.00,0.00}{\textbf{#1}}}
\newcommand{\ExtensionTok}[1]{#1}
\newcommand{\FloatTok}[1]{\textcolor[rgb]{0.00,0.00,0.81}{#1}}
\newcommand{\FunctionTok}[1]{\textcolor[rgb]{0.13,0.29,0.53}{\textbf{#1}}}
\newcommand{\ImportTok}[1]{#1}
\newcommand{\InformationTok}[1]{\textcolor[rgb]{0.56,0.35,0.01}{\textbf{\textit{#1}}}}
\newcommand{\KeywordTok}[1]{\textcolor[rgb]{0.13,0.29,0.53}{\textbf{#1}}}
\newcommand{\NormalTok}[1]{#1}
\newcommand{\OperatorTok}[1]{\textcolor[rgb]{0.81,0.36,0.00}{\textbf{#1}}}
\newcommand{\OtherTok}[1]{\textcolor[rgb]{0.56,0.35,0.01}{#1}}
\newcommand{\PreprocessorTok}[1]{\textcolor[rgb]{0.56,0.35,0.01}{\textit{#1}}}
\newcommand{\RegionMarkerTok}[1]{#1}
\newcommand{\SpecialCharTok}[1]{\textcolor[rgb]{0.81,0.36,0.00}{\textbf{#1}}}
\newcommand{\SpecialStringTok}[1]{\textcolor[rgb]{0.31,0.60,0.02}{#1}}
\newcommand{\StringTok}[1]{\textcolor[rgb]{0.31,0.60,0.02}{#1}}
\newcommand{\VariableTok}[1]{\textcolor[rgb]{0.00,0.00,0.00}{#1}}
\newcommand{\VerbatimStringTok}[1]{\textcolor[rgb]{0.31,0.60,0.02}{#1}}
\newcommand{\WarningTok}[1]{\textcolor[rgb]{0.56,0.35,0.01}{\textbf{\textit{#1}}}}
\usepackage{longtable,booktabs,array}
\usepackage{calc} % for calculating minipage widths
\usepackage{caption}
% Make caption package work with longtable
\makeatletter
\def\fnum@table{\tablename~\thetable}
\makeatother
\setlength{\emergencystretch}{3em} % prevent overfull lines
\providecommand{\tightlist}{%
  \setlength{\itemsep}{0pt}\setlength{\parskip}{0pt}}
\usepackage{bookmark}
\IfFileExists{xurl.sty}{\usepackage{xurl}}{} % add URL line breaks if available
\urlstyle{same}
\hypersetup{
  pdftitle={Regresión Logística},
  pdfauthor={ML Finanzas - Clase 4},
  hidelinks,
  pdfcreator={LaTeX via pandoc}}

\title{Regresión Logística}
\subtitle{Todo lo que necesitas saber}
\author{ML Finanzas - Clase 4}
\date{26/02/2026}

\begin{document}
\frame{\titlepage}

\begin{frame}
\begin{block}{Tabla de contenidos}
\phantomsection\label{tabla-de-contenidos}
\begin{enumerate}
\tightlist
\item
  Matemáticas detrás de la regresión logística
\item
  Implementación desde cero
\item
  Regresión logística regularizada
\item
  Clasificación multiclase: Uno-contra-todos
\item
  Comparación con sklearn y SVM
\end{enumerate}
\end{block}
\end{frame}

\begin{frame}
\begin{block}{¿Qué es la Regresión Logística?}
\phantomsection\label{quuxe9-es-la-regresiuxf3n-loguxedstica}
\begin{itemize}
\tightlist
\item
  Algoritmo de \textbf{Aprendizaje Supervisado}
\item
  Tiene la etiqueta (variable dependiente) para cada ejemplo de
  entrenamiento
\item
  \textbf{Objetivo:} Encontrar la mejor frontera de decisión que separa
  los datos en diferentes clases
\item
  Responde preguntas como: \emph{``¿Es un gato o un perro?''}
\end{itemize}
\end{block}
\end{frame}

\begin{frame}
\begin{block}{Hipótesis y Función Sigmoide}
\phantomsection\label{hipuxf3tesis-y-funciuxf3n-sigmoide}
La hipótesis de la regresión logística se define como:

\[h_\theta(x) = g(\theta^Tx)\]

donde \(g\) es la función sigmoide:

\[g(z) = \frac{1}{1+e^{-z}}\]
\end{block}
\end{frame}

\begin{frame}
\begin{block}{¿Por qué es un modelo lineal?}
\phantomsection\label{por-quuxe9-es-un-modelo-lineal}
\(g\) modela la relación entre \(p(x) = Pr(y=1|x)\) y \(x\).

Se puede demostrar que:

\[\log\left(\frac{p(x)}{1 - p(x)}\right) = \theta^Tx\]

La cantidad \(\frac{p(x)}{1 - p(x)}\) se llama \textbf{odds}. El logit
(log-odds) es \textbf{lineal} en \(x\).
\end{block}
\end{frame}

\begin{frame}
\begin{block}{Función de Costo}
\phantomsection\label{funciuxf3n-de-costo}
No podemos usar mínimos cuadrados (no es convexa). Usamos \textbf{máxima
verosimilitud}:

\[J(\theta) = \frac{1}{m}\sum_{i=1}^m\left(-y_i \log(h_i(\theta)) - (1-y_i)\log(1-h_i(\theta))\right)\]

\begin{itemize}
\tightlist
\item
  Buscamos \(\theta\) que minimice \(J(\theta)\)
\item
  Función de costo convexa para regresión logística
\end{itemize}
\end{block}
\end{frame}

\begin{frame}
\begin{block}{Descenso de Gradiente}
\phantomsection\label{descenso-de-gradiente}
Actualización en cada iteración:

\[\theta_j := \theta_j - \frac{\alpha}{m} \sum_{i=1}^{m} (h_\theta(x^{(i)}) - y^{(i)}) x_j^{(i)}\]

\begin{itemize}
\tightlist
\item
  \(\alpha\): tasa de aprendizaje
\item
  Actualizar simultáneamente todos los \(\theta_j\)
\item
  Con cada paso, los parámetros se acercan al óptimo
\end{itemize}
\end{block}
\end{frame}

\begin{frame}
\begin{block}{Ejemplo: Admisión Universitaria}
\phantomsection\label{ejemplo-admisiuxf3n-universitaria}
\textbf{Problema:} Determinar la probabilidad de admisión de aspirantes
basándose en: - Calificación Examen 1 - Calificación Examen 2

\textbf{Datos:} Histórico de aspirantes con decisión de admisión (sí/no)
\end{block}
\end{frame}

\begin{frame}[fragile]
\begin{block}{Funciones Clave - Implementación}
\phantomsection\label{funciones-clave---implementaciuxf3n}
\begin{Shaded}
\begin{Highlighting}[]
\KeywordTok{def}\NormalTok{ sigmoid(z):}
    \ControlFlowTok{return} \DecValTok{1}\OperatorTok{/}\NormalTok{(}\DecValTok{1}\OperatorTok{+}\NormalTok{np.exp(}\OperatorTok{{-}}\NormalTok{z))}

\KeywordTok{def}\NormalTok{ costFunction(theta, X, y):}
\NormalTok{    h }\OperatorTok{=}\NormalTok{ sigmoid(X }\OperatorTok{@}\NormalTok{ theta)}
    \ControlFlowTok{return}\NormalTok{ (}\OperatorTok{{-}}\NormalTok{y.T }\OperatorTok{@}\NormalTok{ np.log(h) }\OperatorTok{{-}}\NormalTok{ (}\DecValTok{1}\OperatorTok{{-}}\NormalTok{y).T }\OperatorTok{@}\NormalTok{ np.log(}\DecValTok{1}\OperatorTok{{-}}\NormalTok{h)) }\OperatorTok{/}\NormalTok{ m}
\end{Highlighting}
\end{Shaded}
\end{block}
\end{frame}

\begin{frame}[fragile]
\begin{block}{Algoritmo BFGS}
\phantomsection\label{algoritmo-bfgs}
\begin{itemize}
\tightlist
\item
  Alternativa al descenso de gradiente
\item
  Algoritmo \textbf{Broyden--Fletcher--Goldfarb--Shanno}
\item
  Evita quedar atrapado en óptimos locales
\item
  Usado en \texttt{scipy.optimize.fmin\_bfgs}
\end{itemize}
\end{block}
\end{frame}

\begin{frame}
\begin{block}{Regresión Logística Regularizada}
\phantomsection\label{regresiuxf3n-loguxedstica-regularizada}
\textbf{Problema:} Datos que no se separan con una línea recta

\textbf{Solución:} 1. \textbf{Mapeo de características} - términos
polinomiales hasta potencia 6 2. \textbf{Regularización} - penalizar
parámetros grandes con \(\lambda\)
\end{block}
\end{frame}

\begin{frame}
\begin{block}{Función de Costo Regularizada}
\phantomsection\label{funciuxf3n-de-costo-regularizada}
\[J(\theta) = \frac{1}{m}\sum_{i=1}^m\left(-y_i \log(h_i) - (1-y_i)\log(1-h_i)\right) + \frac{\lambda}{2m}\sum_{j=1}^n\theta_j^2\]

\begin{itemize}
\tightlist
\item
  \textbf{No regularizar} \(\theta_0\) (término de bias)
\item
  \(\lambda\) controla el balance sesgo-varianza
\end{itemize}
\end{block}
\end{frame}

\begin{frame}
\begin{block}{Clasificación Multiclase: Uno-contra-todos}
\phantomsection\label{clasificaciuxf3n-multiclase-uno-contra-todos}
Para \(y \in \{0, 1, ..., n\}\):

\begin{itemize}
\tightlist
\item
  Entrenar \(n+1\) clasificadores binarios
\item
  Cada uno predice \(P(y = i | x; \theta)\)
\item
  \textbf{Predicción:} \(\arg\max_i h_\theta^{(i)}(x)\)
\end{itemize}
\end{block}
\end{frame}

\begin{frame}
\begin{block}{Comparación: Implementación vs sklearn}
\phantomsection\label{comparaciuxf3n-implementaciuxf3n-vs-sklearn}
\begin{longtable}[]{@{}ll@{}}
\toprule\noalign{}
Método & Precisión típica \\
\midrule\noalign{}
\endhead
Nuestra implementación (One-vs-all) & \textasciitilde82\% \\
sklearn LogisticRegression & \textasciitilde92\% \\
\bottomrule\noalign{}
\end{longtable}

sklearn usa optimización más eficiente y configuración por defecto
mejor.
\end{block}
\end{frame}

\begin{frame}
\begin{block}{Comparación: Regresión Logística vs SVM}
\phantomsection\label{comparaciuxf3n-regresiuxf3n-loguxedstica-vs-svm}
\begin{longtable}[]{@{}ll@{}}
\toprule\noalign{}
Modelo & Precisión (dígitos MNIST) \\
\midrule\noalign{}
\endhead
Regresión Logística (sklearn) & \textasciitilde92\% \\
SVM (kernel RBF) & \textasciitilde89\% \\
\bottomrule\noalign{}
\end{longtable}

La elección depende del problema: RL es más interpretable, SVM puede
capturar fronteras no lineales complejas.
\end{block}
\end{frame}

\begin{frame}
\begin{block}{Resumen}
\phantomsection\label{resumen}
\begin{itemize}
\tightlist
\item
  \textbf{Regresión logística:} clasificación binaria/multiclase
\item
  \textbf{Función sigmoide} para probabilidades entre 0 y 1
\item
  \textbf{Descenso de gradiente} o \textbf{BFGS} para optimización
\item
  \textbf{Regularización} para evitar sobreajuste
\item
  \textbf{Uno-contra-todos} para múltiples clases
\end{itemize}
\end{block}
\end{frame}

\begin{frame}
\begin{block}{Referencias}
\phantomsection\label{referencias}
\begin{itemize}
\tightlist
\item
  Notebook original:
  \href{https://github.com/pptr3/personal-blog/blob/master/article/logistic-regression/LogisticRegressionArticle.ipynb}{LogisticRegressionArticle}
\item
  Curso Machine Learning - Andrew Ng (Coursera)
\item
  Documentación:
  \href{https://scikit-learn.org/stable/modules/linear_model.html\#logistic-regression}{scikit-learn
  Logistic Regression}
\end{itemize}
\end{block}
\end{frame}

\begin{frame}{¡Gracias!}
\phantomsection\label{gracias}
\textbf{Preguntas}
\end{frame}

\end{document}
